
\documentclass{article}
\usepackage{colortbl}
\usepackage{makecell}
\usepackage{multirow}
\usepackage{supertabular}

\begin{document}

\newcounter{utterance}

\centering \large Interaction Transcript for game `clean\_up', experiment `1\_medium\_3obj\_pl', episode 4 with Qwen3.6{-}35B{-}A3B{-}FP8{-}without{-}reasoning.
\vspace{24pt}

{ \footnotesize  \setcounter{utterance}{1}
\setlength{\tabcolsep}{0pt}
\begin{supertabular}{c@{$\;$}|p{.15\linewidth}@{}p{.15\linewidth}p{.15\linewidth}p{.15\linewidth}p{.15\linewidth}p{.15\linewidth}}
    \# & $\;$A & \multicolumn{4}{c}{Game Master} & $\;\:$B\\
    \hline

    \theutterance \stepcounter{utterance}  
    & & \multicolumn{4}{p{0.6\linewidth}}{
        \cellcolor[rgb]{0.9,0.9,0.9}{
            \makecell[{{p{\linewidth}}}]{
                \texttt{\tiny{[P1$\langle$GM]}}
                \texttt{Jestem mistrzem gry, a Państwo grają w grę kooperacyjną z następującą siatką jako planszą:} \\
\\ 
\texttt{```} \\
\texttt{ 1234567} \\
\texttt{╔═╤══╤══╗} \\
\texttt{║◌│◌◌│◌◌║ 1} \\
\texttt{╟─┤◌◌│◌◌║ 2} \\
\texttt{║◌│◌◌├──╢ 3} \\
\texttt{╟─┘C◌│◌◌║ 4} \\
\texttt{║◌◌◌◌│◌◌║ 5} \\
\texttt{║◌L┌─┼──╢ 6} \\
\texttt{║◌◌│◌│P◌║ 7} \\
\texttt{╚══╧═╧══╝} \\
\\ 
\texttt{```} \\
\\ 
\texttt{* Górna krawędź wyświetla współrzędne x rosnące w prawo, a prawa krawędź współrzędne y rosnące w dół.} \\
\texttt{* Następujące obiekty są losowo umieszczone na siatce: 'C', 'L', 'P'.} \\
\\ 
\\ 
\texttt{Drugi gracz widzi wariant planszy, na której obiekty są umieszczone w różnych losowych miejscach. Nie można zobaczyć planszy drugiego gracza, a on nie może zobaczyć Państwa planszy.} \\
\\ 
\\ 
\texttt{**Cel:**} \\
\\ 
\texttt{Obaj gracze muszą przesuwać obiekty na swoim tle tak, aby identyczne obiekty znalazły się na tych samych współrzędnych. Muszą Państwo komunikować się z drugim graczem, aby uzgodnić wspólną konfigurację docelową.} \\
\\ 
\\ 
\texttt{**Zasady:**} \\
\\ 
\texttt{* W każdej turze można wysłać dokładnie jedno z dwóch poniższych poleceń:} \\
\texttt{ 1. `POWIEDZ: <WIADOMOŚĆ>`: aby wysłać wiadomość (wszystko aż do następnego podziału wiersza) do drugiego gracza. Przekażę ją Państwa partnerowi.} \\
\texttt{2. `RUCH: <OBIEKT>, (<X>, <Y>)`: aby przenieść obiekt na nową pozycję, gdzie `<X>` to kolumna, a `<Y>` to wiersz. Poinformuję Państwa, czy przeniesienie się powiodło, czy nie.} \\
\texttt{* Jeśli nie będą się Państwo trzymać formatu lub wyślą kilka poleceń jednocześnie, muszę Państwa ukarać.} \\
\texttt{* Jeśli obaj gracze zgromadzą więcej niż 8 kar, obaj przegrają grę.} \\
\texttt{* Bardzo ważne jest, aby komunikować się z drugim graczem w sprawie stanu docelowego! *Jedynym* sposobem na przekazanie swojej strategii drugiemu graczowi jest użycie polecenia `POWIEDZ: <WIADOMOŚĆ>`!} \\
\\ 
\\ 
\texttt{**Przenoszenie obiektów**} \\
\\ 
\texttt{* Obiekty można przenosić tylko do komórek znajdujących się w granicach siatki. Komórka docelowa musi być pusta, tzn. musi zawierać tylko symbol '◌'.} \\
\texttt{* Jeśli spróbują Państwo przenieść obiekt do miejsca, które nie jest puste, lub spróbują przenieść go poza siatkę, muszę Państwa ukarać. Otrzymają Państwo kolejną próbę.} \\
\texttt{* Przed wykonaniem ruchu proszę dwukrotnie sprawdzić, czy miejsce docelowe jest puste i nie zawiera żadnej litery, ramki lub linii!} \\
\\ 
\\ 
\texttt{**Koniec gry**} \\
\\ 
\texttt{Jeśli uważają Państwo, że osiągnęli cel wyrównania wszystkich obiektów, mogą Państwo poprosić drugiego gracza o zakończenie gry, wysyłając `POWIEDZ: zakończono?`. Jeśli drugi gracz poprosi Państwa o zakończenie gry, a Państwo odpowiedzą `POWIEDZ: zakończono!`, gra się zakończy.} \\
\\ 
\texttt{Obaj gracze wygrywają, jeśli gra zakończy się w ciągu 12 rund, gdzie jedna runda jest definiowana jako wysłanie ważnego polecenia przez dwóch graczy.} \\
\\ 
\texttt{**Punktacja:**} \\
\\ 
\texttt{Im bliżej siebie znajdują się identyczne obiekty na obu planszach, tym więcej punktów otrzymują Państwo. Kary zmniejszają liczbę punktów. Czy uda się Państwu pobić rekord?} \\
\\ 
\\ 
\texttt{Proszę wysłać wiadomość do drugiego gracza, aby rozpocząć grę!} \\
            }
        }
    }
    & \\ \\

    \theutterance \stepcounter{utterance}  
    & \multicolumn{4}{p{0.6\linewidth}}{
        \cellcolor[rgb]{0.9,0.9,0.9}{
            \makecell[{{p{\linewidth}}}]{
                \texttt{\tiny{[P1$\rangle$GM]}}
                \texttt{POWIEDZ: Witaj! Zaczynamy. Proszę, podaj mi swoje aktualne współrzędne dla obiektów C, L i P. Użyj formatu: C:(x,y), L:(x,y), P:(x,y). Następnie ustalimy wspólny cel.} \\
            }
        }
    }
    & & \\ \\

    \theutterance \stepcounter{utterance}  
    & & & \multicolumn{2}{p{0.3\linewidth}}{
        \cellcolor[rgb]{0.9,0.9,0.9}{
            \makecell[{{p{\linewidth}}}]{
                \texttt{\tiny{[GM$|$GM]}}
                \texttt{POWIEDZ: Witaj! Zaczynamy. Proszę, podaj mi swoje aktualne współrzędne dla obiektów C, L i P. Użyj formatu: C:(x,y), L:(x,y), P:(x,y). Następnie ustalimy wspólny cel.} \\
            }
        }
    }
    & & \\ \\

    \theutterance \stepcounter{utterance}  
    & & \multicolumn{4}{p{0.6\linewidth}}{
        \cellcolor[rgb]{0.9,0.9,0.9}{
            \makecell[{{p{\linewidth}}}]{
                \texttt{\tiny{[GM$\rangle$P2]}}
                \texttt{Jestem mistrzem gry, a Państwo grają w grę kooperacyjną z następującą siatką jako planszą:} \\
\\ 
\texttt{```} \\
\texttt{ 1234567} \\
\texttt{╔═╤══╤══╗} \\
\texttt{║◌│◌◌│◌◌║ 1} \\
\texttt{╟─┤◌◌│◌◌║ 2} \\
\texttt{║◌│L◌├──╢ 3} \\
\texttt{╟─┘◌◌│◌C║ 4} \\
\texttt{║P◌◌◌│◌◌║ 5} \\
\texttt{║◌◌┌─┼──╢ 6} \\
\texttt{║◌◌│◌│◌◌║ 7} \\
\texttt{╚══╧═╧══╝} \\
\\ 
\texttt{```} \\
\\ 
\texttt{* Górna krawędź wyświetla współrzędne x rosnące w prawo, a prawa krawędź współrzędne y rosnące w dół.} \\
\texttt{* Następujące obiekty są losowo umieszczone na siatce: 'C', 'L', 'P'.} \\
\\ 
\\ 
\texttt{Drugi gracz widzi wariant planszy, na której obiekty są umieszczone w różnych losowych miejscach. Nie można zobaczyć planszy drugiego gracza, a on nie może zobaczyć Państwa planszy.} \\
\\ 
\\ 
\texttt{**Cel:**} \\
\\ 
\texttt{Obaj gracze muszą przesuwać obiekty na swoim tle tak, aby identyczne obiekty znalazły się na tych samych współrzędnych. Muszą Państwo komunikować się z drugim graczem, aby uzgodnić wspólną konfigurację docelową.} \\
\\ 
\\ 
\texttt{**Zasady:**} \\
\\ 
\texttt{* W każdej turze można wysłać dokładnie jedno z dwóch poniższych poleceń:} \\
\texttt{ 1. `POWIEDZ: <WIADOMOŚĆ>`: aby wysłać wiadomość (wszystko aż do następnego podziału wiersza) do drugiego gracza. Przekażę ją Państwa partnerowi.} \\
\texttt{2. `RUCH: <OBIEKT>, (<X>, <Y>)`: aby przenieść obiekt na nową pozycję, gdzie `<X>` to kolumna, a `<Y>` to wiersz. Poinformuję Państwa, czy przeniesienie się powiodło, czy nie.} \\
\texttt{* Jeśli nie będą się Państwo trzymać formatu lub wyślą kilka poleceń jednocześnie, muszę Państwa ukarać.} \\
\texttt{* Jeśli obaj gracze zgromadzą więcej niż 8 kar, obaj przegrają grę.} \\
\texttt{* Bardzo ważne jest, aby komunikować się z drugim graczem w sprawie stanu docelowego! *Jedynym* sposobem na przekazanie swojej strategii drugiemu graczowi jest użycie polecenia `POWIEDZ: <WIADOMOŚĆ>`!} \\
\\ 
\\ 
\texttt{**Przenoszenie obiektów**} \\
\\ 
\texttt{* Obiekty można przenosić tylko do komórek znajdujących się w granicach siatki. Komórka docelowa musi być pusta, tzn. musi zawierać tylko symbol '◌'.} \\
\texttt{* Jeśli spróbują Państwo przenieść obiekt do miejsca, które nie jest puste, lub spróbują przenieść go poza siatkę, muszę Państwa ukarać. Otrzymają Państwo kolejną próbę.} \\
\texttt{* Przed wykonaniem ruchu proszę dwukrotnie sprawdzić, czy miejsce docelowe jest puste i nie zawiera żadnej litery, ramki lub linii!} \\
\\ 
\\ 
\texttt{**Koniec gry**} \\
\\ 
\texttt{Jeśli uważają Państwo, że osiągnęli cel wyrównania wszystkich obiektów, mogą Państwo poprosić drugiego gracza o zakończenie gry, wysyłając `POWIEDZ: zakończono?`. Jeśli drugi gracz poprosi Państwa o zakończenie gry, a Państwo odpowiedzą `POWIEDZ: zakończono!`, gra się zakończy.} \\
\\ 
\texttt{Obaj gracze wygrywają, jeśli gra zakończy się w ciągu 12 rund, gdzie jedna runda jest definiowana jako wysłanie ważnego polecenia przez dwóch graczy.} \\
\\ 
\texttt{**Punktacja:**} \\
\\ 
\texttt{Im bliżej siebie znajdują się identyczne obiekty na obu planszach, tym więcej punktów otrzymują Państwo. Kary zmniejszają liczbę punktów. Czy uda się Państwu pobić rekord?} \\
\\ 
\texttt{Drugi gracz rozpoczął grę wysyłając następującą wiadomość:} \\
\\ 
\texttt{"Witaj! Zaczynamy. Proszę, podaj mi swoje aktualne współrzędne dla obiektów C, L i P. Użyj formatu: C:(x,y), L:(x,y), P:(x,y). Następnie ustalimy wspólny cel."} \\
\\ 
\texttt{Jakie jest twoje pierwsze polecenie?} \\
            }
        }
    }
    & \\ \\

    \theutterance \stepcounter{utterance}  
    & & & \multicolumn{4}{p{0.6\linewidth}}{
        \cellcolor[rgb]{0.9,0.9,0.9}{
            \makecell[{{p{\linewidth}}}]{
                \texttt{\tiny{[GM$\langle$P2]}}
                \texttt{POWIEDZ: Witaj! Oto moje aktualne współrzędne: C:(6,4), L:(3,3), P:(1,5). Proszę, podaj mi swoje współrzędne, abyśmy mogli ustalić wspólny cel.} \\
            }
        }
    }
     \\ \\

    \theutterance \stepcounter{utterance}  
    & & & \multicolumn{2}{p{0.3\linewidth}}{
        \cellcolor[rgb]{0.9,0.9,0.9}{
            \makecell[{{p{\linewidth}}}]{
                \texttt{\tiny{[GM$|$GM]}}
                \texttt{POWIEDZ: Witaj! Oto moje aktualne współrzędne: C:(6,4), L:(3,3), P:(1,5). Proszę, podaj mi swoje współrzędne, abyśmy mogli ustalić wspólny cel.} \\
            }
        }
    }
    & & \\ \\

    \theutterance \stepcounter{utterance}  
    & & \multicolumn{4}{p{0.6\linewidth}}{
        \cellcolor[rgb]{0.9,0.9,0.9}{
            \makecell[{{p{\linewidth}}}]{
                \texttt{\tiny{[P1$\langle$GM]}}
                \texttt{Twoja wiadomość została przekazana drugiemu graczowi.} \\
\\ 
\\ 
\texttt{Obecnie grasz rundę 1 z maksymalnie 12.} \\
\texttt{Łącznie zgromadziliście 0 z 8 kar.} \\
\texttt{Drugi gracz wysłał tę wiadomość:} \\
\\ 
\\ 
\texttt{"Witaj! Oto moje aktualne współrzędne: C:(6,4), L:(3,3), P:(1,5). Proszę, podaj mi swoje współrzędne, abyśmy mogli ustalić wspólny cel."} \\
\\ 
\\ 
\texttt{Jakie jest twoje następne polecenie?} \\
            }
        }
    }
    & \\ \\

    \theutterance \stepcounter{utterance}  
    & \multicolumn{4}{p{0.6\linewidth}}{
        \cellcolor[rgb]{0.9,0.9,0.9}{
            \makecell[{{p{\linewidth}}}]{
                \texttt{\tiny{[P1$\rangle$GM]}}
                \texttt{POWIEDZ: Dziękuję! Oto moje współrzędne: C:(4,4), L:(2,6), P:(7,7). Proponuję jako cel wspólny pozycję C:(6,4), L:(3,3), P:(1,5), czyli Twoje aktualne pozycje. Czy się zgadzasz? Jeśli tak, proszę potwierdź, a następnie wykonaj ruchy do tych pozycji. Ja również wykonam ruchy do tych pozycji.} \\
            }
        }
    }
    & & \\ \\

    \theutterance \stepcounter{utterance}  
    & & & \multicolumn{2}{p{0.3\linewidth}}{
        \cellcolor[rgb]{0.9,0.9,0.9}{
            \makecell[{{p{\linewidth}}}]{
                \texttt{\tiny{[GM$|$GM]}}
                \texttt{POWIEDZ: Dziękuję! Oto moje współrzędne: C:(4,4), L:(2,6), P:(7,7). Proponuję jako cel wspólny pozycję C:(6,4), L:(3,3), P:(1,5), czyli Twoje aktualne pozycje. Czy się zgadzasz? Jeśli tak, proszę potwierdź, a następnie wykonaj ruchy do tych pozycji. Ja również wykonam ruchy do tych pozycji.} \\
            }
        }
    }
    & & \\ \\

    \theutterance \stepcounter{utterance}  
    & & \multicolumn{4}{p{0.6\linewidth}}{
        \cellcolor[rgb]{0.9,0.9,0.9}{
            \makecell[{{p{\linewidth}}}]{
                \texttt{\tiny{[GM$\rangle$P2]}}
                \texttt{Twoja wiadomość została przekazana drugiemu graczowi.} \\
\\ 
\\ 
\texttt{Obecnie grasz rundę 2 z maksymalnie 12.} \\
\texttt{Łącznie zgromadziliście 0 z 8 kar.} \\
\texttt{Drugi gracz wysłał tę wiadomość:} \\
\\ 
\\ 
\texttt{"Dziękuję! Oto moje współrzędne: C:(4,4), L:(2,6), P:(7,7). Proponuję jako cel wspólny pozycję C:(6,4), L:(3,3), P:(1,5), czyli Twoje aktualne pozycje. Czy się zgadzasz? Jeśli tak, proszę potwierdź, a następnie wykonaj ruchy do tych pozycji. Ja również wykonam ruchy do tych pozycji."} \\
\\ 
\\ 
\texttt{Jakie jest twoje następne polecenie?} \\
            }
        }
    }
    & \\ \\

    \theutterance \stepcounter{utterance}  
    & & & \multicolumn{4}{p{0.6\linewidth}}{
        \cellcolor[rgb]{0.9,0.9,0.9}{
            \makecell[{{p{\linewidth}}}]{
                \texttt{\tiny{[GM$\langle$P2]}}
                \texttt{POWIEDZ: Zgadzam się. Cel to C:(6,4), L:(3,3), P:(1,5). Wykonuję ruchy.} \\
            }
        }
    }
     \\ \\

    \theutterance \stepcounter{utterance}  
    & & & \multicolumn{2}{p{0.3\linewidth}}{
        \cellcolor[rgb]{0.9,0.9,0.9}{
            \makecell[{{p{\linewidth}}}]{
                \texttt{\tiny{[GM$|$GM]}}
                \texttt{POWIEDZ: Zgadzam się. Cel to C:(6,4), L:(3,3), P:(1,5). Wykonuję ruchy.} \\
            }
        }
    }
    & & \\ \\

    \theutterance \stepcounter{utterance}  
    & & \multicolumn{4}{p{0.6\linewidth}}{
        \cellcolor[rgb]{0.9,0.9,0.9}{
            \makecell[{{p{\linewidth}}}]{
                \texttt{\tiny{[P1$\langle$GM]}}
                \texttt{Twoja wiadomość została przekazana drugiemu graczowi.} \\
\\ 
\\ 
\texttt{Obecnie grasz rundę 2 z maksymalnie 12.} \\
\texttt{Łącznie zgromadziliście 0 z 8 kar.} \\
\texttt{Drugi gracz wysłał tę wiadomość:} \\
\\ 
\\ 
\texttt{"Zgadzam się. Cel to C:(6,4), L:(3,3), P:(1,5). Wykonuję ruchy."} \\
\\ 
\\ 
\texttt{Jakie jest twoje następne polecenie?} \\
            }
        }
    }
    & \\ \\

    \theutterance \stepcounter{utterance}  
    & \multicolumn{4}{p{0.6\linewidth}}{
        \cellcolor[rgb]{0.9,0.9,0.9}{
            \makecell[{{p{\linewidth}}}]{
                \texttt{\tiny{[P1$\rangle$GM]}}
                \texttt{RUCH: C, (6, 4)} \\
            }
        }
    }
    & & \\ \\

    \theutterance \stepcounter{utterance}  
    & & & \multicolumn{2}{p{0.3\linewidth}}{
        \cellcolor[rgb]{0.9,0.9,0.9}{
            \makecell[{{p{\linewidth}}}]{
                \texttt{\tiny{[GM$|$GM]}}
                \texttt{RUCH: C, (6, 4)} \\
            }
        }
    }
    & & \\ \\

    \theutterance \stepcounter{utterance}  
    & & & \multicolumn{2}{p{0.3\linewidth}}{
        \cellcolor[rgb]{0.9,0.9,0.9}{
            \makecell[{{p{\linewidth}}}]{
                \texttt{\tiny{[GM$|$GM]}}
                \texttt{C moved to (6, 4)} \\
            }
        }
    }
    & & \\ \\

    \theutterance \stepcounter{utterance}  
    & & \multicolumn{4}{p{0.6\linewidth}}{
        \cellcolor[rgb]{0.9,0.9,0.9}{
            \makecell[{{p{\linewidth}}}]{
                \texttt{\tiny{[GM$\rangle$P2]}}
                \texttt{Twoja wiadomość została przekazana drugiemu graczowi.} \\
\\ 
\\ 
\texttt{Obecnie grasz rundę 3 z maksymalnie 12.} \\
\texttt{Łącznie zgromadziliście 0 z 8 kar.} \\
\texttt{Drugi gracz przesunął obiekt na swojej siatce.} \\
\\ 
\\ 
\texttt{Jakie jest twoje następne polecenie?} \\
            }
        }
    }
    & \\ \\

    \theutterance \stepcounter{utterance}  
    & & & \multicolumn{4}{p{0.6\linewidth}}{
        \cellcolor[rgb]{0.9,0.9,0.9}{
            \makecell[{{p{\linewidth}}}]{
                \texttt{\tiny{[GM$\langle$P2]}}
                \texttt{RUCH: P, (1,5)} \\
            }
        }
    }
     \\ \\

    \theutterance \stepcounter{utterance}  
    & & & \multicolumn{2}{p{0.3\linewidth}}{
        \cellcolor[rgb]{0.9,0.9,0.9}{
            \makecell[{{p{\linewidth}}}]{
                \texttt{\tiny{[GM$|$GM]}}
                \texttt{RUCH: P, (1,5)} \\
            }
        }
    }
    & & \\ \\

    \theutterance \stepcounter{utterance}  
    & & & \multicolumn{2}{p{0.3\linewidth}}{
        \cellcolor[rgb]{0.9,0.9,0.9}{
            \makecell[{{p{\linewidth}}}]{
                \texttt{\tiny{[GM$|$GM]}}
                \texttt{Kara: (1,5) nie jest puste, ale zawiera 'P'.} \\
\\ 
\texttt{Łącznie zgromadziliście 1 z 8 kar.} \\
\texttt{Spróbuj ponownie!} \\
            }
        }
    }
    & & \\ \\

    \theutterance \stepcounter{utterance}  
    & & \multicolumn{4}{p{0.6\linewidth}}{
        \cellcolor[rgb]{0.9,0.9,0.9}{
            \makecell[{{p{\linewidth}}}]{
                \texttt{\tiny{[GM$\rangle$P2]}}
                \texttt{Kara: (1,5) nie jest puste, ale zawiera 'P'.} \\
\\ 
\texttt{Łącznie zgromadziliście 1 z 8 kar.} \\
\texttt{Spróbuj ponownie!} \\
            }
        }
    }
    & \\ \\

    \theutterance \stepcounter{utterance}  
    & & & \multicolumn{4}{p{0.6\linewidth}}{
        \cellcolor[rgb]{0.9,0.9,0.9}{
            \makecell[{{p{\linewidth}}}]{
                \texttt{\tiny{[GM$\langle$P2]}}
                \texttt{RUCH: P, (1,1)} \\
            }
        }
    }
     \\ \\

    \theutterance \stepcounter{utterance}  
    & & & \multicolumn{2}{p{0.3\linewidth}}{
        \cellcolor[rgb]{0.9,0.9,0.9}{
            \makecell[{{p{\linewidth}}}]{
                \texttt{\tiny{[GM$|$GM]}}
                \texttt{RUCH: P, (1,1)} \\
            }
        }
    }
    & & \\ \\

    \theutterance \stepcounter{utterance}  
    & & & \multicolumn{2}{p{0.3\linewidth}}{
        \cellcolor[rgb]{0.9,0.9,0.9}{
            \makecell[{{p{\linewidth}}}]{
                \texttt{\tiny{[GM$|$GM]}}
                \texttt{P moved to (1, 1)} \\
            }
        }
    }
    & & \\ \\

    \theutterance \stepcounter{utterance}  
    & & \multicolumn{4}{p{0.6\linewidth}}{
        \cellcolor[rgb]{0.9,0.9,0.9}{
            \makecell[{{p{\linewidth}}}]{
                \texttt{\tiny{[P1$\langle$GM]}}
                \texttt{Pomyślnie przeniesiono 'C' do (6,4). Twoja zaktualizowana siatka wygląda następująco:} \\
\texttt{```} \\
\texttt{ 1234567} \\
\texttt{╔═╤══╤══╗} \\
\texttt{║◌│◌◌│◌◌║ 1} \\
\texttt{╟─┤◌◌│◌◌║ 2} \\
\texttt{║◌│◌◌├──╢ 3} \\
\texttt{╟─┘◌◌│C◌║ 4} \\
\texttt{║◌◌◌◌│◌◌║ 5} \\
\texttt{║◌L┌─┼──╢ 6} \\
\texttt{║◌◌│◌│P◌║ 7} \\
\texttt{╚══╧═╧══╝} \\
\texttt{```} \\
\texttt{Obecnie grasz rundę 3 z maksymalnie 12.} \\
\texttt{Łącznie zgromadziliście 1 z 8 kar.} \\
\texttt{Drugi gracz przesunął obiekt na swojej siatce.} \\
\\ 
\\ 
\texttt{Jakie jest twoje następne polecenie?} \\
            }
        }
    }
    & \\ \\

    \theutterance \stepcounter{utterance}  
    & \multicolumn{4}{p{0.6\linewidth}}{
        \cellcolor[rgb]{0.9,0.9,0.9}{
            \makecell[{{p{\linewidth}}}]{
                \texttt{\tiny{[P1$\rangle$GM]}}
                \texttt{RUCH: L, (3, 3)} \\
            }
        }
    }
    & & \\ \\

    \theutterance \stepcounter{utterance}  
    & & & \multicolumn{2}{p{0.3\linewidth}}{
        \cellcolor[rgb]{0.9,0.9,0.9}{
            \makecell[{{p{\linewidth}}}]{
                \texttt{\tiny{[GM$|$GM]}}
                \texttt{RUCH: L, (3, 3)} \\
            }
        }
    }
    & & \\ \\

    \theutterance \stepcounter{utterance}  
    & & & \multicolumn{2}{p{0.3\linewidth}}{
        \cellcolor[rgb]{0.9,0.9,0.9}{
            \makecell[{{p{\linewidth}}}]{
                \texttt{\tiny{[GM$|$GM]}}
                \texttt{L moved to (3, 3)} \\
            }
        }
    }
    & & \\ \\

    \theutterance \stepcounter{utterance}  
    & & \multicolumn{4}{p{0.6\linewidth}}{
        \cellcolor[rgb]{0.9,0.9,0.9}{
            \makecell[{{p{\linewidth}}}]{
                \texttt{\tiny{[GM$\rangle$P2]}}
                \texttt{Pomyślnie przeniesiono 'P' do (1,1). Twoja zaktualizowana siatka wygląda następująco:} \\
\texttt{```} \\
\texttt{ 1234567} \\
\texttt{╔═╤══╤══╗} \\
\texttt{║P│◌◌│◌◌║ 1} \\
\texttt{╟─┤◌◌│◌◌║ 2} \\
\texttt{║◌│L◌├──╢ 3} \\
\texttt{╟─┘◌◌│◌C║ 4} \\
\texttt{║◌◌◌◌│◌◌║ 5} \\
\texttt{║◌◌┌─┼──╢ 6} \\
\texttt{║◌◌│◌│◌◌║ 7} \\
\texttt{╚══╧═╧══╝} \\
\texttt{```} \\
\texttt{Obecnie grasz rundę 4 z maksymalnie 12.} \\
\texttt{Łącznie zgromadziliście 1 z 8 kar.} \\
\texttt{Drugi gracz przesunął obiekt na swojej siatce.} \\
\\ 
\\ 
\texttt{Jakie jest twoje następne polecenie?} \\
            }
        }
    }
    & \\ \\

    \theutterance \stepcounter{utterance}  
    & & & \multicolumn{4}{p{0.6\linewidth}}{
        \cellcolor[rgb]{0.9,0.9,0.9}{
            \makecell[{{p{\linewidth}}}]{
                \texttt{\tiny{[GM$\langle$P2]}}
                \texttt{RUCH: L, (3,3)} \\
            }
        }
    }
     \\ \\

    \theutterance \stepcounter{utterance}  
    & & & \multicolumn{2}{p{0.3\linewidth}}{
        \cellcolor[rgb]{0.9,0.9,0.9}{
            \makecell[{{p{\linewidth}}}]{
                \texttt{\tiny{[GM$|$GM]}}
                \texttt{RUCH: L, (3,3)} \\
            }
        }
    }
    & & \\ \\

    \theutterance \stepcounter{utterance}  
    & & & \multicolumn{2}{p{0.3\linewidth}}{
        \cellcolor[rgb]{0.9,0.9,0.9}{
            \makecell[{{p{\linewidth}}}]{
                \texttt{\tiny{[GM$|$GM]}}
                \texttt{Kara: (3,3) nie jest puste, ale zawiera 'L'.} \\
\\ 
\texttt{Łącznie zgromadziliście 2 z 8 kar.} \\
\texttt{Spróbuj ponownie!} \\
            }
        }
    }
    & & \\ \\

    \theutterance \stepcounter{utterance}  
    & & \multicolumn{4}{p{0.6\linewidth}}{
        \cellcolor[rgb]{0.9,0.9,0.9}{
            \makecell[{{p{\linewidth}}}]{
                \texttt{\tiny{[GM$\rangle$P2]}}
                \texttt{Kara: (3,3) nie jest puste, ale zawiera 'L'.} \\
\\ 
\texttt{Łącznie zgromadziliście 2 z 8 kar.} \\
\texttt{Spróbuj ponownie!} \\
            }
        }
    }
    & \\ \\

    \theutterance \stepcounter{utterance}  
    & & & \multicolumn{4}{p{0.6\linewidth}}{
        \cellcolor[rgb]{0.9,0.9,0.9}{
            \makecell[{{p{\linewidth}}}]{
                \texttt{\tiny{[GM$\langle$P2]}}
                \texttt{RUCH: L, (3,1)} \\
            }
        }
    }
     \\ \\

    \theutterance \stepcounter{utterance}  
    & & & \multicolumn{2}{p{0.3\linewidth}}{
        \cellcolor[rgb]{0.9,0.9,0.9}{
            \makecell[{{p{\linewidth}}}]{
                \texttt{\tiny{[GM$|$GM]}}
                \texttt{RUCH: L, (3,1)} \\
            }
        }
    }
    & & \\ \\

    \theutterance \stepcounter{utterance}  
    & & & \multicolumn{2}{p{0.3\linewidth}}{
        \cellcolor[rgb]{0.9,0.9,0.9}{
            \makecell[{{p{\linewidth}}}]{
                \texttt{\tiny{[GM$|$GM]}}
                \texttt{L moved to (3, 1)} \\
            }
        }
    }
    & & \\ \\

    \theutterance \stepcounter{utterance}  
    & & \multicolumn{4}{p{0.6\linewidth}}{
        \cellcolor[rgb]{0.9,0.9,0.9}{
            \makecell[{{p{\linewidth}}}]{
                \texttt{\tiny{[P1$\langle$GM]}}
                \texttt{Pomyślnie przeniesiono 'L' do (3,3). Twoja zaktualizowana siatka wygląda następująco:} \\
\texttt{```} \\
\texttt{ 1234567} \\
\texttt{╔═╤══╤══╗} \\
\texttt{║◌│◌◌│◌◌║ 1} \\
\texttt{╟─┤◌◌│◌◌║ 2} \\
\texttt{║◌│L◌├──╢ 3} \\
\texttt{╟─┘◌◌│C◌║ 4} \\
\texttt{║◌◌◌◌│◌◌║ 5} \\
\texttt{║◌◌┌─┼──╢ 6} \\
\texttt{║◌◌│◌│P◌║ 7} \\
\texttt{╚══╧═╧══╝} \\
\texttt{```} \\
\texttt{Obecnie grasz rundę 4 z maksymalnie 12.} \\
\texttt{Łącznie zgromadziliście 2 z 8 kar.} \\
\texttt{Drugi gracz przesunął obiekt na swojej siatce.} \\
\\ 
\\ 
\texttt{Jakie jest twoje następne polecenie?} \\
            }
        }
    }
    & \\ \\

    \theutterance \stepcounter{utterance}  
    & \multicolumn{4}{p{0.6\linewidth}}{
        \cellcolor[rgb]{0.9,0.9,0.9}{
            \makecell[{{p{\linewidth}}}]{
                \texttt{\tiny{[P1$\rangle$GM]}}
                \texttt{RUCH: P, (1, 5)} \\
            }
        }
    }
    & & \\ \\

    \theutterance \stepcounter{utterance}  
    & & & \multicolumn{2}{p{0.3\linewidth}}{
        \cellcolor[rgb]{0.9,0.9,0.9}{
            \makecell[{{p{\linewidth}}}]{
                \texttt{\tiny{[GM$|$GM]}}
                \texttt{RUCH: P, (1, 5)} \\
            }
        }
    }
    & & \\ \\

    \theutterance \stepcounter{utterance}  
    & & & \multicolumn{2}{p{0.3\linewidth}}{
        \cellcolor[rgb]{0.9,0.9,0.9}{
            \makecell[{{p{\linewidth}}}]{
                \texttt{\tiny{[GM$|$GM]}}
                \texttt{P moved to (1, 5)} \\
            }
        }
    }
    & & \\ \\

    \theutterance \stepcounter{utterance}  
    & & \multicolumn{4}{p{0.6\linewidth}}{
        \cellcolor[rgb]{0.9,0.9,0.9}{
            \makecell[{{p{\linewidth}}}]{
                \texttt{\tiny{[GM$\rangle$P2]}}
                \texttt{Pomyślnie przeniesiono 'L' do (3,1). Twoja zaktualizowana siatka wygląda następująco:} \\
\texttt{```} \\
\texttt{ 1234567} \\
\texttt{╔═╤══╤══╗} \\
\texttt{║P│L◌│◌◌║ 1} \\
\texttt{╟─┤◌◌│◌◌║ 2} \\
\texttt{║◌│◌◌├──╢ 3} \\
\texttt{╟─┘◌◌│◌C║ 4} \\
\texttt{║◌◌◌◌│◌◌║ 5} \\
\texttt{║◌◌┌─┼──╢ 6} \\
\texttt{║◌◌│◌│◌◌║ 7} \\
\texttt{╚══╧═╧══╝} \\
\texttt{```} \\
\texttt{Obecnie grasz rundę 5 z maksymalnie 12.} \\
\texttt{Łącznie zgromadziliście 2 z 8 kar.} \\
\texttt{Drugi gracz przesunął obiekt na swojej siatce.} \\
\\ 
\\ 
\texttt{Jakie jest twoje następne polecenie?} \\
            }
        }
    }
    & \\ \\

    \theutterance \stepcounter{utterance}  
    & & & \multicolumn{4}{p{0.6\linewidth}}{
        \cellcolor[rgb]{0.9,0.9,0.9}{
            \makecell[{{p{\linewidth}}}]{
                \texttt{\tiny{[GM$\langle$P2]}}
                \texttt{RUCH: C, (6,4)} \\
            }
        }
    }
     \\ \\

    \theutterance \stepcounter{utterance}  
    & & & \multicolumn{2}{p{0.3\linewidth}}{
        \cellcolor[rgb]{0.9,0.9,0.9}{
            \makecell[{{p{\linewidth}}}]{
                \texttt{\tiny{[GM$|$GM]}}
                \texttt{RUCH: C, (6,4)} \\
            }
        }
    }
    & & \\ \\

    \theutterance \stepcounter{utterance}  
    & & & \multicolumn{2}{p{0.3\linewidth}}{
        \cellcolor[rgb]{0.9,0.9,0.9}{
            \makecell[{{p{\linewidth}}}]{
                \texttt{\tiny{[GM$|$GM]}}
                \texttt{C moved to (6, 4)} \\
            }
        }
    }
    & & \\ \\

    \theutterance \stepcounter{utterance}  
    & & \multicolumn{4}{p{0.6\linewidth}}{
        \cellcolor[rgb]{0.9,0.9,0.9}{
            \makecell[{{p{\linewidth}}}]{
                \texttt{\tiny{[P1$\langle$GM]}}
                \texttt{Pomyślnie przeniesiono 'P' do (1,5). Twoja zaktualizowana siatka wygląda następująco:} \\
\texttt{```} \\
\texttt{ 1234567} \\
\texttt{╔═╤══╤══╗} \\
\texttt{║◌│◌◌│◌◌║ 1} \\
\texttt{╟─┤◌◌│◌◌║ 2} \\
\texttt{║◌│L◌├──╢ 3} \\
\texttt{╟─┘◌◌│C◌║ 4} \\
\texttt{║P◌◌◌│◌◌║ 5} \\
\texttt{║◌◌┌─┼──╢ 6} \\
\texttt{║◌◌│◌│◌◌║ 7} \\
\texttt{╚══╧═╧══╝} \\
\texttt{```} \\
\texttt{Obecnie grasz rundę 5 z maksymalnie 12.} \\
\texttt{Łącznie zgromadziliście 2 z 8 kar.} \\
\texttt{Drugi gracz przesunął obiekt na swojej siatce.} \\
\\ 
\\ 
\texttt{Jakie jest twoje następne polecenie?} \\
            }
        }
    }
    & \\ \\

    \theutterance \stepcounter{utterance}  
    & \multicolumn{4}{p{0.6\linewidth}}{
        \cellcolor[rgb]{0.9,0.9,0.9}{
            \makecell[{{p{\linewidth}}}]{
                \texttt{\tiny{[P1$\rangle$GM]}}
                \texttt{POWIEDZ: Wszystkie obiekty są na swoich miejscach. Czy Ty też je wszystkie przeniosłeś? Jeśli tak, odpowiedz "tak", a następnie wyślij "zakończono?".} \\
            }
        }
    }
    & & \\ \\

    \theutterance \stepcounter{utterance}  
    & & & \multicolumn{2}{p{0.3\linewidth}}{
        \cellcolor[rgb]{0.9,0.9,0.9}{
            \makecell[{{p{\linewidth}}}]{
                \texttt{\tiny{[GM$|$GM]}}
                \texttt{POWIEDZ: Wszystkie obiekty są na swoich miejscach. Czy Ty też je wszystkie przeniosłeś? Jeśli tak, odpowiedz "tak", a następnie wyślij "zakończono?".} \\
            }
        }
    }
    & & \\ \\

    \theutterance \stepcounter{utterance}  
    & & \multicolumn{4}{p{0.6\linewidth}}{
        \cellcolor[rgb]{0.9,0.9,0.9}{
            \makecell[{{p{\linewidth}}}]{
                \texttt{\tiny{[GM$\rangle$P2]}}
                \texttt{Pomyślnie przeniesiono 'C' do (6,4). Twoja zaktualizowana siatka wygląda następująco:} \\
\texttt{```} \\
\texttt{ 1234567} \\
\texttt{╔═╤══╤══╗} \\
\texttt{║P│L◌│◌◌║ 1} \\
\texttt{╟─┤◌◌│◌◌║ 2} \\
\texttt{║◌│◌◌├──╢ 3} \\
\texttt{╟─┘◌◌│C◌║ 4} \\
\texttt{║◌◌◌◌│◌◌║ 5} \\
\texttt{║◌◌┌─┼──╢ 6} \\
\texttt{║◌◌│◌│◌◌║ 7} \\
\texttt{╚══╧═╧══╝} \\
\texttt{```} \\
\texttt{Obecnie grasz rundę 6 z maksymalnie 12.} \\
\texttt{Łącznie zgromadziliście 2 z 8 kar.} \\
\texttt{Drugi gracz wysłał tę wiadomość:} \\
\\ 
\\ 
\texttt{"Wszystkie obiekty są na swoich miejscach. Czy Ty też je wszystkie przeniosłeś? Jeśli tak, odpowiedz "tak", a następnie wyślij "zakończono?"."} \\
\\ 
\\ 
\texttt{Jakie jest twoje następne polecenie?} \\
            }
        }
    }
    & \\ \\

    \theutterance \stepcounter{utterance}  
    & & & \multicolumn{4}{p{0.6\linewidth}}{
        \cellcolor[rgb]{0.9,0.9,0.9}{
            \makecell[{{p{\linewidth}}}]{
                \texttt{\tiny{[GM$\langle$P2]}}
                \texttt{POWIEDZ: tak. POWIEDZ: zakończono?} \\
            }
        }
    }
     \\ \\

    \theutterance \stepcounter{utterance}  
    & & & \multicolumn{2}{p{0.3\linewidth}}{
        \cellcolor[rgb]{0.9,0.9,0.9}{
            \makecell[{{p{\linewidth}}}]{
                \texttt{\tiny{[GM$|$GM]}}
                \texttt{POWIEDZ: tak. POWIEDZ: zakończono?} \\
            }
        }
    }
    & & \\ \\

    \theutterance \stepcounter{utterance}  
    & & & \multicolumn{2}{p{0.3\linewidth}}{
        \cellcolor[rgb]{0.9,0.9,0.9}{
            \makecell[{{p{\linewidth}}}]{
                \texttt{\tiny{[GM$|$GM]}}
                \texttt{Invalid format: head is not empty: 'POWIEDZ: tak. '} \\
            }
        }
    }
    & & \\ \\

    \theutterance \stepcounter{utterance}  
    & & \multicolumn{4}{p{0.6\linewidth}}{
        \cellcolor[rgb]{0.9,0.9,0.9}{
            \makecell[{{p{\linewidth}}}]{
                \texttt{\tiny{[GM$\rangle$P2]}}
                \texttt{Kara: Twoja wiadomość nie może zawierać niczego przed poleceniem!} \\
\\ 
\texttt{Upewnij się, że twoja odpowiedź zawiera tylko `POWIEDZ: <WIADOMOŚĆ>` lub `RUCH: <OBIEKT>, (<X>, <Y>)` i nic więcej!} \\
\\ 
\texttt{Łącznie zgromadziliście 3 z 8 kar.} \\
\texttt{Spróbuj ponownie!} \\
            }
        }
    }
    & \\ \\

    \theutterance \stepcounter{utterance}  
    & & & \multicolumn{4}{p{0.6\linewidth}}{
        \cellcolor[rgb]{0.9,0.9,0.9}{
            \makecell[{{p{\linewidth}}}]{
                \texttt{\tiny{[GM$\langle$P2]}}
                \texttt{POWIEDZ: tak. zakończono?} \\
            }
        }
    }
     \\ \\

    \theutterance \stepcounter{utterance}  
    & & & \multicolumn{2}{p{0.3\linewidth}}{
        \cellcolor[rgb]{0.9,0.9,0.9}{
            \makecell[{{p{\linewidth}}}]{
                \texttt{\tiny{[GM$|$GM]}}
                \texttt{POWIEDZ: tak. zakończono?} \\
            }
        }
    }
    & & \\ \\

    \theutterance \stepcounter{utterance}  
    & & \multicolumn{4}{p{0.6\linewidth}}{
        \cellcolor[rgb]{0.9,0.9,0.9}{
            \makecell[{{p{\linewidth}}}]{
                \texttt{\tiny{[P1$\langle$GM]}}
                \texttt{Twoja wiadomość została przekazana drugiemu graczowi.} \\
\\ 
\\ 
\texttt{Obecnie grasz rundę 6 z maksymalnie 12.} \\
\texttt{Łącznie zgromadziliście 3 z 8 kar.} \\
\texttt{Drugi gracz wysłał tę wiadomość:} \\
\\ 
\\ 
\texttt{"tak. zakończono?"} \\
\\ 
\\ 
\texttt{Jakie jest twoje następne polecenie?} \\
            }
        }
    }
    & \\ \\

    \theutterance \stepcounter{utterance}  
    & \multicolumn{4}{p{0.6\linewidth}}{
        \cellcolor[rgb]{0.9,0.9,0.9}{
            \makecell[{{p{\linewidth}}}]{
                \texttt{\tiny{[P1$\rangle$GM]}}
                \texttt{POWIEDZ: zakończono!} \\
            }
        }
    }
    & & \\ \\

    \theutterance \stepcounter{utterance}  
    & & & \multicolumn{2}{p{0.3\linewidth}}{
        \cellcolor[rgb]{0.9,0.9,0.9}{
            \makecell[{{p{\linewidth}}}]{
                \texttt{\tiny{[GM$|$GM]}}
                \texttt{POWIEDZ: zakończono!} \\
            }
        }
    }
    & & \\ \\

    \theutterance \stepcounter{utterance}  
    & & & \multicolumn{2}{p{0.3\linewidth}}{
        \cellcolor[rgb]{0.9,0.9,0.9}{
            \makecell[{{p{\linewidth}}}]{
                \texttt{\tiny{[GM$|$GM]}}
                \texttt{true} \\
            }
        }
    }
    & & \\ \\

    \theutterance \stepcounter{utterance}  
    & & & \multicolumn{2}{p{0.3\linewidth}}{
        \cellcolor[rgb]{0.9,0.9,0.9}{
            \makecell[{{p{\linewidth}}}]{
                \texttt{\tiny{[GM$|$GM]}}
                \texttt{Initial states:} \\
\texttt{```} \\
\texttt{Player 1:} \\
\texttt{ 1234567} \\
\texttt{╔═╤══╤══╗} \\
\texttt{║◌│◌◌│◌◌║ 1} \\
\texttt{╟─┤◌◌│◌◌║ 2} \\
\texttt{║◌│◌◌├──╢ 3} \\
\texttt{╟─┘C◌│◌◌║ 4} \\
\texttt{║◌◌◌◌│◌◌║ 5} \\
\texttt{║◌L┌─┼──╢ 6} \\
\texttt{║◌◌│◌│P◌║ 7} \\
\texttt{╚══╧═╧══╝} \\
\\ 
\\ 
\texttt{Player 2:} \\
\texttt{ 1234567} \\
\texttt{╔═╤══╤══╗} \\
\texttt{║◌│◌◌│◌◌║ 1} \\
\texttt{╟─┤◌◌│◌◌║ 2} \\
\texttt{║◌│L◌├──╢ 3} \\
\texttt{╟─┘◌◌│◌C║ 4} \\
\texttt{║P◌◌◌│◌◌║ 5} \\
\texttt{║◌◌┌─┼──╢ 6} \\
\texttt{║◌◌│◌│◌◌║ 7} \\
\texttt{╚══╧═╧══╝} \\
\\ 
\texttt{```} \\
            }
        }
    }
    & & \\ \\

    \theutterance \stepcounter{utterance}  
    & & & \multicolumn{2}{p{0.3\linewidth}}{
        \cellcolor[rgb]{0.9,0.9,0.9}{
            \makecell[{{p{\linewidth}}}]{
                \texttt{\tiny{[GM$|$GM]}}
                \texttt{End states:} \\
\texttt{```} \\
\texttt{Player 1:} \\
\texttt{ 1234567} \\
\texttt{╔═╤══╤══╗} \\
\texttt{║◌│◌◌│◌◌║ 1} \\
\texttt{╟─┤◌◌│◌◌║ 2} \\
\texttt{║◌│L◌├──╢ 3} \\
\texttt{╟─┘◌◌│C◌║ 4} \\
\texttt{║P◌◌◌│◌◌║ 5} \\
\texttt{║◌◌┌─┼──╢ 6} \\
\texttt{║◌◌│◌│◌◌║ 7} \\
\texttt{╚══╧═╧══╝} \\
\\ 
\\ 
\texttt{Player 2:} \\
\texttt{ 1234567} \\
\texttt{╔═╤══╤══╗} \\
\texttt{║P│L◌│◌◌║ 1} \\
\texttt{╟─┤◌◌│◌◌║ 2} \\
\texttt{║◌│◌◌├──╢ 3} \\
\texttt{╟─┘◌◌│C◌║ 4} \\
\texttt{║◌◌◌◌│◌◌║ 5} \\
\texttt{║◌◌┌─┼──╢ 6} \\
\texttt{║◌◌│◌│◌◌║ 7} \\
\texttt{╚══╧═╧══╝} \\
\\ 
\texttt{```} \\
            }
        }
    }
    & & \\ \\

    \theutterance \stepcounter{utterance}  
    & & & \multicolumn{2}{p{0.3\linewidth}}{
        \cellcolor[rgb]{0.9,0.9,0.9}{
            \makecell[{{p{\linewidth}}}]{
                \texttt{\tiny{[GM$|$GM]}}
                \texttt{* success: True} \\
\texttt{* lose: False} \\
\texttt{* aborted: False} \\
\texttt{{-}{-}{-}{-}{-}{-}{-}} \\
\texttt{* Shifts: 4.00} \\
\texttt{* Max Shifts: 4.00} \\
\texttt{* Min Shifts: 2.00} \\
\texttt{* End Distance Sum: 6.00} \\
\texttt{* Init Distance Sum: 12.55} \\
\texttt{* Expected Distance Sum: 12.57} \\
\texttt{* Penalties: 3.00} \\
\texttt{* Max Penalties: 8.00} \\
\texttt{* Rounds: 6.00} \\
\texttt{* Max Rounds: 12.00} \\
\texttt{* Object Count: 3.00} \\
            }
        }
    }
    & & \\ \\

\end{supertabular}
}

\end{document}
